% spec body template — the slug-echo body filename (<slug>.tex, per #295)
% lives at <thread>.{N}/<slug>.tex. This is a normative specification: every
% claim is true of the implementation (code_ref) or explicitly marked
% target-state. A spec's body is LaTeX (SKILL.md §Output format) — multi-file
% friendly via \input{sections/*.tex}. Delete these comments before authoring.
% Replace bracketed placeholders.
%
% ADOPTION NOTE: if you are adopting an existing multi-file LaTeX spec, you do
% NOT start from this template — place the existing tree per SKILL.md
% §"Adopting an existing spec". This skeleton is only for the deferred
% draft-from-scratch / stub-scaffold case (spec-draft step 4 "Scaffold").

\documentclass[11pt]{article}

\usepackage[margin=1in]{geometry}
\usepackage{graphicx}
\usepackage{amsmath,amssymb}
\usepackage{booktabs}
\usepackage{hyperref}

\title{[System] Specification}
\author{[Maintainer / team]}
\date{\today}

\begin{document}
\maketitle
\tableofcontents

\section{Scope and conformance}
% What this spec normatively governs. State the RFC 2119 keyword discipline
% (MUST / SHALL / SHOULD / MAY) once here and use it consistently. Name the
% implementation this spec describes (the code_ref companion) so a reader
% knows the source of truth the normative claims are checked against.

\section{Definitions and notation}
% Define EVERY term used in a normative predicate below. An undefined term
% used in a MUST/SHALL is the "Undefined normative term" critical flag.

\section{[Component / message / structure]}
% The normative content. For each constant, struct/field layout, formula, and
% validity predicate: state it unambiguously and falsifiably (dim 3), and keep
% it consistent with every other section that states the same quantity (dim 2).
% Each claim either matches the implementation (code_ref) or is explicitly
% marked target-state with the gap tracked.

\subsection{Constants}
% A grep-able constants table makes internal-consistency and (Phase 3)
% deterministic constant-extraction checks cheap. State each constant once,
% authoritatively; reference it elsewhere rather than restating a literal.
\begin{center}
\begin{tabular}{@{}lll@{}}
\toprule
Name & Value & Notes \\
\midrule
[CONST\_A] & [value] & [what it governs] \\
\bottomrule
\end{tabular}
\end{center}

\section{Validity predicates}
% The formal block/message/state-validity predicates. Every symbol used here
% must be defined in §Definitions. A predicate must be satisfiable (dim 5).

\section{Implementation status}
% Live vs. target per component. Where the implementation and the spec
% diverge, mark it HERE rather than silently rewriting either side. NOTE: the
% first-class implementation-status register (and its cross-check in the audit)
% is Phase 2 scope (#707) — this section is the placeholder it formalizes.

\section{Revision history}
% Versioning discipline (dim 7): a reader must know which version they hold and
% what changed between versions.
\begin{center}
\begin{tabular}{@{}lll@{}}
\toprule
Version & Date & Summary of normative changes \\
\midrule
1 & \today & Initial draft \\
\bottomrule
\end{tabular}
\end{center}

\end{document}
